\chapter{동기식 설계: 클럭, 레지스터, 조합 논리}\label{ch:sync}

앞 장의 회로는 입력이 바뀌면 출력이 따라 바뀔 뿐 아무것도 기억하지 못했다. 쓸모 있는 회로는 상태를
기억해야 한다. 카운터는 지금까지 센 수를, 봄베는 676가닥의 통전 여부를 기억해야 한다. 이 장은 디지털
회로가 기억과 시간을 다루는 표준적인 방법을 설명한다.

\section{두 종류의 부품}

디지털 회로는 두 종류의 부품으로 이루어진다.

\begin{itemize}
\item \term{조합 논리}{combinational logic}: 현재 입력만으로 출력이 정해지는 회로. 게이트를 엮은 것이며
기억이 없다. 입력이 바뀌면 \emph{얼마 뒤에} 출력이 바뀐다. 이 ``얼마''가 지연이다.
\item \term{레지스터}{register}: 값을 기억하는 부품. 플립플롭(flip-flop)이라는 1비트 기억 소자를 묶은
것이다. \term{클럭}{clock}이라는 주기적인 신호가 0에서 1로 올라가는 순간(상승 에지)에만 입력을 받아들여
출력으로 내보내고, 다음 상승 에지까지 그 값을 붙들고 있다.
\end{itemize}

\Why 왜 값을 클럭 에지에서만 바꾸는가. 조합 논리의 출력은 입력이 바뀐 직후에는 믿을 수 없다. 신호가
게이트를 지나는 데 시간이 걸리고 경로마다 그 시간이 달라서, 잠깐 동안 엉뚱한 값이 나타났다 사라진다.
회로 전체가 ``이제 값이 안정되었으니 받아 적자''고 합의할 기준 시각이 필요한데, 그것이 클럭 에지다.
모든 레지스터가 같은 클럭에 맞춰 값을 바꾸도록 설계하는 방식을 \term{동기식 설계}{synchronous design}라
하고, FPGA 설계는 사실상 전부 이 방식이다.

\General 그림~\ref{fig:g-dff}는 플립플롭 하나의 동작이다. 입력 \code{d}는 아무 때나 바뀌어도 되지만
출력 \code{q}는 상승 에지에서만 바뀐다. \code{d}에 잠깐 나타났다 사라진 값(두 에지 사이의 짧은 펄스)은
\code{q}에 전혀 나타나지 않는다.

\begin{figure}[ht]
\centering
\begin{tikzpicture}[x=1cm,y=1cm]
  \node[reg,minimum width=14mm,minimum height=16mm] (ff) at (0,0) {};
  \node[font=\ttfamily\small,anchor=west] at ([xshift=1pt]ff.west) {d};
  \node[font=\ttfamily\small,anchor=east] at ([xshift=-1pt]ff.east) {q};
  \pic at ([xshift=0mm]ff.south) {clkmark};
  \draw[<-] (ff.west) -- ++(-0.9,0);
  \draw[->] (ff.east) -- ++(0.9,0);
  \draw[<-] (ff.south) -- ++(0,-0.5) node[below,font=\ttfamily\small] {clk};
\end{tikzpicture}
\hspace{8mm}
\begin{tikztimingtable}[timing/slope=0.1,timing/coldist=4pt,xscale=1.5,yscale=1.25,semithick]
  \code{clk} & 8{.5L .5H} \\
  \code{d}   & .2L 1.5H .5L .3H 1.2L 2H 2.3L \\
  \code{q}   & .5L 1H 2L 2H 2.5L \\
\extracode
  \begin{pgfonlayer}{background}
    \vertlines[cgray!50,dashed]{0.5,1.5,...,7.5}
  \end{pgfonlayer}
\end{tikztimingtable}
\caption{플립플롭. 왼쪽은 기호, 오른쪽은 동작이다. \code{q}는 상승 에지(점선)에서의 \code{d}를 받아 다음
에지까지 유지한다.}
\label{fig:g-dff}
\end{figure}

\section{레지스터 사이의 한 클럭}

동기식 회로의 모든 부분은 그림~\ref{fig:g-rcr}의 꼴로 볼 수 있다. 레지스터에서 나온 값이 조합 논리를
지나 다음 레지스터(같은 레지스터일 수도 있다)의 입력에 도착한다.

\begin{figure}[ht]
\centering
\begin{tikzpicture}[x=1cm,y=1cm]
  \node[reg,minimum height=14mm] (r1) at (0,0) {레지\\스터};
  \node[cloud,minimum width=44mm,minimum height=14mm] (c) at (4.2,0) {조합 논리\\{\footnotesize 지연 $t_{\mathrm{logic}}$}};
  \node[reg,minimum height=14mm] (r2) at (8.4,0) {레지\\스터};
  \draw[->] (r1) -- (c); \draw[->] (c) -- (r2);
  \draw[<-] (r1.south) -- ++(0,-0.6) coordinate (k1);
  \draw[<-] (r2.south) -- ++(0,-0.6) coordinate (k2);
  \draw (k1) -- (k2) node[midway,below,font=\ttfamily\small] {clk};
  \pic at (r1.south) {clkmark}; \pic at (r2.south) {clkmark};
  \draw[<->,cred] (0.6,1.15) -- node[above,note,text=cred] {이 구간의 지연이 클럭 주기보다 짧아야 한다} (7.8,1.15);
\end{tikzpicture}
\caption{동기식 회로의 기본 단위. 에지에서 출발한 값이 다음 에지 전에 도착해야 한다.}
\label{fig:g-rcr}
\end{figure}

\Why 이 꼴에서 동기식 설계의 유일한 시간 규칙이 나온다. \emph{한 에지에서 레지스터를 떠난 값은 조합
논리를 지나 다음 에지 전에 다음 레지스터에 도착해 있어야 한다.} 늦으면 다음 레지스터가 덜 계산된 값을
받아 적는다. 따라서 가장 느린 조합 논리 경로의 지연이 클럭 주기의 하한을, 곧 회로가 동작할 수 있는
최고 클럭 주파수를 정한다. 이 규칙만 지키면 조합 논리 안에서 값이 어떻게 출렁이든 상관이 없다.
설계자가 신경 쓸 시간의 단위가 ``나노초''에서 ``클럭 몇 번''으로 바뀌는 것이 동기식 설계의 이점이다.

\General 가장 단순한 예는 카운터다(그림~\ref{fig:g-counter}). 레지스터의 출력에 1을 더하는 조합 논리를
달고, 그 결과를 다시 레지스터의 입력에 넣는다. 에지마다 값이 1씩 커진다.

\begin{lstlisting}[language=SV,numbers=none]
logic [3:0] count;
always_ff @(posedge clk)
  if (reset) count <= 4'd0;
  else       count <= count + 4'd1;
\end{lstlisting}

\begin{figure}[ht]
\centering
\begin{tikzpicture}[x=1cm,y=1cm]
  \node[reg,minimum height=13mm,minimum width=13mm] (r) at (0,0) {\code{count}};
  \node[cloud,minimum width=16mm] (a) at (2.6,0) {$+1$};
  \draw[->] (r) -- (a);
  \draw[->] (a.east) -- ++(0.5,0) |- (1.3,1.2) -| ([xshift=-7mm]r.west) -- (r.west);
  \pic at (r.south) {clkmark};
  \draw[<-] (r.south) -- ++(0,-0.45) node[below,font=\ttfamily\small] {clk};
\end{tikzpicture}
\hspace{6mm}
\begin{tikztimingtable}[timing/slope=0.1,timing/coldist=4pt,xscale=1.9,yscale=1.25,semithick]
  \code{clk}       & 5{.5L .5H} \\
  \code{count}     & .5D{0} 1D{1} 1D{2} 1D{3} 1D{4} .5D{5} \\
  \code{count+1}   & .5D{1} 1D{2} 1D{3} 1D{4} 1D{5} .5D{6} \\
\extracode
  \begin{pgfonlayer}{background}
    \vertlines[cgray!50,dashed]{0.5,1.5,...,4.5}
  \end{pgfonlayer}
\end{tikztimingtable}
\caption{카운터. 조합 논리는 ``다음 값''을 늘 계산하고 있고, 레지스터는 에지마다 그것을 받아 ``현재 값''으로
삼는다.}
\label{fig:g-counter}
\end{figure}

카운터에서 익혀 둘 것은 이름 두 개다. 레지스터에 담긴 \emph{현재 값}과, 조합 논리가 계산해 둔 \emph{다음
값}. 동기식 회로를 읽는 것은 결국 ``다음 값이 현재 값에서 어떻게 계산되는가''를 읽는 일이다.

\Ours 봄베의 기준 RTL은 이 꼴 그대로다. \code{live\_q}가 676비트의 현재 상태를 담은 레지스터이고,
\code{live\_n}이 조합 논리(코드~\ref{lst:g-spread})가 계산한 다음 상태다. 이름 끝의 \code{\_q}와
\code{\_n}은 각각 레지스터의 출력과 다음 값(next)을 뜻하는 관례다. 카운터의 ``$+1$'' 자리에
``한 걸음 퍼뜨리기''가 들어갔을 뿐이다(그림~\ref{fig:g-liveloop}).

\begin{figure}[ht]
\centering
\begin{tikzpicture}[x=1cm,y=1cm]
  \node[reg,minimum height=16mm,minimum width=16mm] (r) at (0,0) {\code{live\_q}\\{\footnotesize 676비트}};
  \node[cloud,minimum width=44mm,minimum height=16mm] (f) at (4.6,0) {한 걸음 퍼뜨리기\\{\footnotesize 스크램블러 + 대각선 보드}};
  \draw[->] (r) -- (f);
  \draw[->] (f.east) -- ++(0.6,0) |- (2.3,1.5) node[pos=0.75,above,font=\footnotesize] {\code{live\_n}} -| ([xshift=-8mm]r.west) -- (r.west);
  \pic at (r.south) {clkmark};
  \draw[<-] (r.south) -- ++(0,-0.45) node[below,font=\ttfamily\small] {clk};
\end{tikzpicture}
\caption{\code{bombe\_core}의 뼈대. 그림~\ref{fig:g-counter}의 카운터와 같은 꼴이다.}
\label{fig:g-liveloop}
\end{figure}

실제 봄베에서는 전류가 전선을 따라 저절로, 사실상 즉시 퍼진다. 그것을 게이트로 그대로 옮기면 출력이
자기 입력으로 돌아오는 고리가 되는데, 그런 고리는 한번 1이 되면 꺼지지 않고 설계 도구가 지연을 계산할
수도 없다. 그래서 고리를 레지스터로 끊고 클럭마다 한 걸음씩 퍼뜨린다. 이 설계 결정은 \ref{ch:rtl-core}장에서
자세히 다룬다.

\section{\texorpdfstring{\code{<=}와 \code{=}}{<=와 =}: 동시에 갱신한다는 것}

SystemVerilog에는 대입 기호가 둘 있다. 레지스터를 기술하는 \code{always\_ff} 안에서는 \code{<=}(논블로킹
대입)를, 조합 논리를 기술하는 \code{always\_comb} 안에서는 \code{=}(블로킹 대입)를 쓴다.

\Why 클럭 에지에서 모든 레지스터는 \emph{동시에} 값을 바꾼다. 그러므로 어떤 레지스터의 새 값을 계산할 때
쓰이는 다른 레지스터의 값은 모두 \emph{에지 직전의} 값이어야 한다. \code{<=}는 바로 그 뜻이다. ``오른쪽을
지금의 값으로 계산해 두었다가, 에지에서 한꺼번에 왼쪽에 넣는다.'' 반면 \code{=}는 조합 논리를 적을 때
``이 전선은 이 식의 결과다''를 차례로 이어 적기 위한 것이다(코드~\ref{lst:g-scr}의 \code{x}처럼).

\General 두 레지스터의 값을 맞바꾸는 회로를 보자.

\begin{lstlisting}[language=SV,numbers=none]
always_ff @(posedge clk) begin
  a <= b;
  b <= a;
end
\end{lstlisting}

프로그램이라면 첫 줄에서 \code{a}가 덮어써지므로 임시 변수가 필요하다. 회로에서는 필요 없다. 두 줄 모두
에지 직전의 값을 읽고, 에지에서 동시에 쓴다. 그림으로는 두 레지스터의 출력이 서로의 입력에 이어져 있을
뿐이다(그림~\ref{fig:g-swap}).

\begin{figure}[ht]
\centering
\begin{tikzpicture}[x=1cm,y=1cm]
  \node[reg,minimum width=11mm] (a) at (0,0) {\code{a}};
  \node[reg,minimum width=11mm] (b) at (3,0) {\code{b}};
  \draw[->] (a.east) -- ++(0.4,0) |- (1.5,0.9) -| ([xshift=-4mm]b.west) -- (b.west);
  \draw[->] (b.east) -- ++(0.5,0) |- (1.5,-1.0) -| ([xshift=-5mm]a.west) -- (a.west);
\end{tikzpicture}
\hspace{10mm}
\begin{tikztimingtable}[timing/slope=0.1,timing/coldist=4pt,xscale=1.8,yscale=1.25,semithick]
  \code{clk} & 5{.5L .5H} \\
  \code{a}  & .5D{3} 1D{8} 1D{3} 1D{8} 1D{3} .5D{8} \\
  \code{b}   & .5D{8} 1D{3} 1D{8} 1D{3} 1D{8} .5D{3} \\
\extracode
  \begin{pgfonlayer}{background}
    \vertlines[cgray!50,dashed]{0.5,1.5,...,4.5}
  \end{pgfonlayer}
\end{tikztimingtable}
\caption{맞바꾸기. 두 레지스터가 같은 에지에서 서로의 옛 값을 받는다.}
\label{fig:g-swap}
\end{figure}

\Ours \code{bombe\_core}에서 상태를 갱신하는 부분은 다음과 같다. 세 줄의 \code{<=}는 차례로 일어나는
일이 아니라 같은 에지에서 한꺼번에 일어나는 일이다. \code{iterations + 16'd1}에 쓰인 \code{iterations}도,
조건식의 \code{live\_q}도 모두 에지 직전의 값이다.

\begin{lstlisting}[language=SV,numbers=none]
SPREAD: begin
  live_q     <= live_n;
  iterations <= iterations + 16'd1;
  if (live_n == live_q || (early_exit && test_n == FULL)) begin
    state <= IDLE;
    done  <= 1'b1;
  end
end
\end{lstlisting}

\section{``차례로''는 상태 기계로 만든다}\label{sec:fsm}

\Why 회로는 모든 것이 동시에 동작하므로, ``먼저 이것을 하고 그다음에 저것을 한다''는 순서는 저절로
생기지 않는다. 순서가 필요하면 ``지금 몇 번째 단계인가''를 기억하는 레지스터를 두고, 단계에 따라 다른
동작을 하게 한다. 이것이 \term{상태 기계}{finite state machine, FSM}다. 프로그램에서는 처리 장치가
``지금 몇 번째 줄인가''를 대신 기억해 주지만, 회로에서는 설계자가 그 기억을 직접 만든다.

\General 버튼을 누르면 LED를 정해진 시간 동안 켜 두는 회로를 생각하자. 상태는 둘이다. 기다리는
\code{IDLE}과 켜 두는 \code{ON}. \code{IDLE}에서 버튼이 눌리면 타이머를 0으로 하고 \code{ON}으로 간다.
\code{ON}에서는 클럭마다 타이머를 올리다가 정해진 값이 되면 \code{IDLE}로 돌아간다
(그림~\ref{fig:g-fsm} 왼쪽).

\Ours \code{bombe\_core}의 상태 기계도 상태가 둘이고 구조가 같다(그림~\ref{fig:g-fsm} 오른쪽).
\code{IDLE}에서 시작 신호가 오면 상태를 초기화하고 \code{SPREAD}로 간다. \code{SPREAD}에서는 클럭마다 한
걸음씩 퍼뜨리다가 더 변하지 않으면 완료를 알리고 \code{IDLE}로 돌아간다. 타이머 대신 ``변화가 없는가''가
끝내는 조건일 뿐이다.

\begin{figure}[ht]
\centering
\begin{tikzpicture}[x=1cm,y=1cm,
   st/.style={draw,circle,minimum size=14mm,fill=cgreenbg,draw=cgreen,font=\small\ttfamily},
   sb/.style={draw,circle,minimum size=14mm,fill=cbluebg,draw=cblue,font=\small\ttfamily}]
  \node[font=\small\bfseries] at (1.8,2.4) {일반적인 예: LED 타이머};
  \node[st] (i) at (0,0) {IDLE};
  \node[st] (o) at (3.6,0) {ON};
  \draw[->] (i) to[bend left=25] node[above,note] {버튼 눌림\\타이머 $\leftarrow 0$} (o);
  \draw[->] (o) to[bend left=25] node[below,note] {타이머 = 끝값} (i);
  \draw[->] (o) to[out=30,in=-30,looseness=5] node[right,note] {타이머 $+1$} (o);
  \begin{scope}[xshift=8.6cm]
  \node[font=\small\bfseries] at (1.8,2.4) {\code{bombe\_core}};
  \node[sb] (i2) at (0,0) {IDLE};
  \node[sb] (s2) at (3.6,0) {SPREAD};
  \draw[->] (i2) to[bend left=25] node[above,note] {\code{start}\\상태 초기화} (s2);
  \draw[->] (s2) to[bend left=25] node[below,note] {변화 없음: \code{done}} (i2);
  \draw[->] (s2) to[out=30,in=-30,looseness=5] node[right,note] {한 걸음} (s2);
  \end{scope}
\end{tikzpicture}
\caption{상태 기계 두 개. 상태를 기억하는 것도 레지스터이고, 화살표는 ``다음 상태''를 정하는 조합 논리다.}
\label{fig:g-fsm}
\end{figure}

\section{정리}

\begin{itemize}
\item 회로는 조합 논리(계산)와 레지스터(기억)로 이루어진다.
\item 레지스터는 클럭 에지에서만 값을 바꾸고, 모든 레지스터가 동시에 바꾼다.
\item 설계자가 적는 것은 ``현재 값에서 다음 값이 어떻게 나오는가''다.
\item 레지스터 사이의 조합 논리는 한 클럭 주기 안에 끝나야 한다. 이 제약이 최고 속도를 정한다.
\item 순서가 필요하면 상태 기계를 만든다.
\end{itemize}
